% =============================================================================
\section{Discussion}\label{sec:discussion}
% =============================================================================
The completed factorial experiment does not support a universal advantage for complex-valued networks. Most complex configurations do not exceed the fixed magnitude baseline, while a subset does. The central conclusion is therefore conditional: the usefulness of complex-valued processing depends on the coordinate representation and the architectural operators used. This heterogeneity is precisely why the study evaluates a parameter-matched factorial grid rather than a single real-versus-complex pair.

Input representation provides the largest broad contrast in the observed grid. Image-domain configurations had mean AUC 0.7102, compared with 0.6377 for Fourier-domain configurations. In the 240 matched architectural pairs, all other factorial choices are held fixed within each pair, and the image-domain result was higher in 196/240 pairs (81.7\%), with mean paired image-minus-Fourier $\Delta$AUC of +0.0725 and median $\Delta$AUC of +0.0852. These are descriptive properties of the finite factorial grid, not causal or population-level statistical effects. A direct Fourier-domain classifier may impose a less favorable locality and aggregation bias for this slice-level anatomical task, but this does not imply that Fourier-domain information is intrinsically inferior. The Fourier arm is an FFT-derived representation of the prepared, phase-aligned complex images rather than untouched scanner raw k-space.

Normalization shows a visible but context-dependent contrast. Complex BatchNorm has the higher marginal mean AUC in the observed grid (0.685 versus 0.663 for RMSNorm). Its covariance whitening and learned Cartesian affine transformation may interact more favorably with this network and training setup than the study-specific ComplexRMSNorm2d implementation. The former is a literature-established complex normalization family, whereas the latter is an implementation choice designed here around running complex power. The observed marginal difference therefore does not establish universal superiority for Complex BatchNorm.

Standard versus widely-linear convolution shows a comparatively modest marginal contrast (mean AUC 0.677 versus 0.671), and the activation means are also close relative to the domain separation. This does not make widely-linear processing or any activation family useless. The value of conjugate flexibility and of the nonlinearities depends on their interaction with representation domain, normalization, pooling, and stream architecture. The factorial results should consequently not be read as a universal ranking of convolution or activation operators.

The combined stream/interaction factor is likewise domain dependent. It is one implemented five-level factor---complex-only/none, complex-only/holographic, dual/none, dual/modulus gate, and dual/holographic---rather than two independent stream and interaction dimensions. For example, the image-domain dual/none level has mean AUC 0.741, whereas the corresponding Fourier-domain level has mean AUC 0.617. This contrast illustrates that an interaction or stream design cannot be assigned a domain-independent value from the present grid. The dual-stream architectures and their optional mechanisms should be interpreted together with the coordinate system in which they operate.

The canonical pooling comparison provides a clean local architectural check because the three complex configurations hold image-domain input, ComplexRMSNorm2d, standard complex convolution, a complex-only stream, no interaction, and modReLU fixed, varying only pooling. The canonical max-pool reached AUC approximately 0.7527, compared with approximately 0.7283 for canonical median-pool and 0.7156 for canonical average-pool; their AP values were 0.148, 0.100, and 0.112, respectively, versus 0.106 for the real baseline. Max pooling therefore performs best among these three canonical choices, but this local comparison does not establish universal superiority across the full factorial grid. The comparison also connects the observed pooling behavior to the phase-coherence terms in Eq.~\eqref{eq:complex-average-coherence} without isolating pooling from the broader grid context.

Average precision is important here because the test cohort contains approximately 5\% positive slices. Across the joint AUC--AP landscape, 114/480 complex configurations (23.8\%) exceeded the real baseline in both metrics; the corresponding counts were 101/240 (42.1\%) for image-domain configurations and 13/240 (5.4\%) for Fourier-domain configurations. AUC and AP were strongly related over the grid, with Pearson $r=0.755$ and Spearman $\rho=0.788$, but they are not interchangeable under this class imbalance: AUC summarizes ranking across false-positive rates, whereas AP is sensitive to the precision achieved among relatively rare positives. These counts and correlations describe the finite factorial grid and do not constitute significance tests or population inference.

The maximum observed AUC of 0.8162 remains useful as an exploratory description of the architecture surface, not as the principal claim of the study. All 480 configurations were evaluated on the same fixed held-out test cohort, and selecting the highest observed test AUC from that grid introduces model-selection optimism. Accordingly, 0.8162 is not an unbiased estimate of the generalization performance of a selected final architecture.

Several limitations bound the interpretation. Each reported architecture has one training realization at fixed seed 10383, so stochastic optimization variability is not quantified. The same test cohort, containing only 63 positive slices, is reused across the 480 complex evaluations; this supports a finite-grid comparison but not independent uncertainty estimates for architecture selection. Metrics are reported at the slice level even though slices within a patient are correlated. The task is a slice-level PI-RADS 3--5 versus PI-RADS 1--2 classification, where PI-RADS 3--5 denotes radiological suspicion rather than histopathologically proven clinically significant cancer, and it is not a patient-level diagnosis.

The real comparator is magnitude-only, so the experiment does not separately identify the value of retaining real-imaginary information from the value of complex algebra. An information-matched real-valued network receiving concatenated real and imaginary channels would separate these mechanisms more cleanly. The study also uses a prepared four-coil T2-weighted subset and an FFT-derived Fourier representation rather than untouched scanner raw k-space. T2-only scope does not represent a complete biparametric or multiparametric clinical prostate MRI assessment, where complementary diffusion information and patient-level aggregation can matter.

The factorial experiment therefore establishes a controlled characterization of where complex-valued neural processing helps or fails under matched scalar capacity, rather than a universal claim of complex-valued superiority. The usefulness of complex-valued modeling for prostate T2 MRI classification is conditional on the coordinate representation and the architectural operators used; confirming a selected architecture requires repeated training realizations, patient-level analysis, and independent validation.
